\documentclass[10pt,openany,a4paper]{article}
\usepackage{tikz, pgfplots, tkz-euclide,calc}
\usetikzlibrary{patterns,shapes.arrows,calc}
\usepackage{fancyhdr}
\usepackage{enumerate,enumitem}
\usepackage{cancel}
\usepackage{varwidth}

% TAMBAHKAN PACKAGE SENDIRI KALAU KURANG

\usepackage{geometry}
\geometry{
	left = 22mm,
	right = 22mm,
	top = 25mm,
	bottom = 25mm,
}
\usepackage{graphicx}
\usepackage{subcaption}
\usepackage{soul}
\usepackage{hyperref}
\usepackage{float}
\usepackage{setspace}
\usepackage{afterpage}
\usepackage{xcolor}
\usepackage[nosectionbib]{apacite}
\usepackage{lmodern}
\usepackage{mathtools}
\usepackage{upgreek}
\usepackage{siunitx}
\usepackage{physics}
\usepackage{hyperref}
\renewcommand{\baselinestretch}{1.3}
\usepackage{multicol}
\usepackage{hyperref,ragged2e,amsmath,multicol,gensymb,setspace,
fancyhdr,amsfonts,tikz,pgfplots,nccmath,enumerate,verbatim,amsthm,physics}
\usepgfplotslibrary{polar,fillbetween}
% Disable tikz externalization to avoid requiring -shell-escape
\usepgfplotslibrary{external}
\tikzexternalize[prefix=tikz/]
\tikzset{reuse path/.code={\pgfsyssoftpath@setcurrentpath{#1}}}
\usetikzlibrary{shapes.geometric,arrows}
\newcommand\mylog[1]{\mathop{{}^{#1}\mathrm{log}}}
\pgfplotsset{compat=1.15}
\pgfplotsset{my style/.append style={axis x line=middle, axis y line=
middle, xlabel={$x$}, ylabel={$y$}, axis equal }}
\usepackage{etoolbox}
% remove stray \makeatother
\newtheorem{defn}{Definisi}[section]
\newtheorem{teo}[defn]{Teorema}
\newtheorem{thm}{Teorema}[section]
\newtheorem{lemma}{Lemma}
\newtheorem{lemmas}[defn]{Lemma}
\newtheorem{cor}[defn]{Akibat}
\newtheorem{asumsi}[defn]{Asumsi}
\theoremstyle{definition}
\newtheorem{con}{Contoh}[section]
\theoremstyle{definition}
\theoremstyle{plain}
\newtheorem{prop}{Proposisi}[section]
\renewcommand{\proofname}{Bukti}
\renewcommand{\thethm}{\arabic{chapter}.\arabic{thm}}

\newcommand{\normlp}[1]{\left\|#1\right\|_{L^p_{\lambda}(0,T;H)}}% Fungsi norm (||x||)
\newcommand{\R}{\mathbb{R}}
\renewcommand{\natural}{\mathbb{N}}
\newcommand{\Xpl}{X^R_{p,\lambda}}
\newcommand{\Xpls}{X^R_{p,\lambda^*}}
\newcommand{\qpl}{q^R_{p,\lambda}}
\newcommand{\qpls}{q^R_{p,\lambda^*}}
\pagestyle{fancy}
\fancyhf{}
\lhead{Halaman \thepage}
\rhead{Quiz 2 Kalkulus 2 Kelas 29\\ (\href{https://instagram.com/ahmadzakiyudin_/}{@ahmadzakiyudin\_})}
\hypersetup{
    colorlinks=true,
    linkcolor=blue,
    filecolor=blue,      
    urlcolor=blue,
}
\newenvironment{jawab}[2][tes]{
% \fbox{\begin{minipage}{14.65cm}
\textbf{Kata kunci materi:} #1.\\
\textbf{Penyelesaian:} #2
% \end{minipage}}
}{}

\newenvironment{jawab2}[2][tes]{
\fbox{\begin{minipage}{14.65cm}
 #2
\end{minipage}}}{}

\newenvironment{jawaban}[2][tes]{
\colorlet{oldcolor}{.}
        \color{red}
        \setlength{\fboxrule}{1pt}\boxed{\color{oldcolor}#2}\color{oldcolor}.
}{}

\newenvironment{jawaban2}[2][tes]{
\colorlet{oldcolor}{.}
        \color{red}
        \setlength{\fboxrule}{1pt}\fbox{\color{oldcolor}#2}\color{oldcolor}.
}{}

\newenvironment{penting}[2][tes]{
\colorlet{oldcolor}{.}
\color{blue}
\fbox{\color{oldcolor}#2}\color{oldcolor}
}{}

\begin{document}
\setlength\fboxsep{3pt}
% ensure enough header height for fancyhdr
\setlength{\headheight}{24pt}
\section{Questions}
\begin{enumerate}
    \item Given the curve $y=\sqrt{4-x^2}$, $-1\leq x \leq 1$, determine:
    \begin{enumerate}
        \item The area under the curve on the given interval bounded by the $x$-axis.
        \item If the area in (a) is revolved about the $x$-axis, determine the resulting volume.
        \item Determine the centroid of the solid in (b) without using Pappus' theorem.
        \item If the interval $-1\leq x \leq 1$ is removed, determine the arc length of the curve in first quadrant, without using integration.
        \item Based on the arc obtained in (d), if it is revolved about the $y$-axis, determine the centroid of the resulting surface area using Pappus' theorem.
    \end{enumerate}
    \item Using integration, determine the arc length of the curve $y=5x$ on the interval $0\leq x\leq 1$.
\end{enumerate}
\section{Solutions}
\begin{enumerate}
    \item \begin{enumerate}
        \item The area under the curve can be calculated using the formula
        \begin{align*}
            A &= \int_{-1}^{1} \sqrt{4-x^2} \, dx \\
        \end{align*}
        Let $x=2\sin\theta$, then $dx=2\cos\theta d\theta$ and the limits of integration change to $\theta=-\pi/6$ and $\theta=\pi/6$. Thus,
        \begin{align*}
            A &= \int_{-\pi/6}^{\pi/6} \sqrt{4-4\sin^2\theta} \cdot 2\cos\theta \, d\theta \\
            &= \int_{-\pi/6}^{\pi/6} 4\cos^2\theta \, d\theta \\
            &= 4 \int_{-\pi/6}^{\pi/6} \cos^2\theta \, d\theta \\
            &= 4 \int_{-\pi/6}^{\pi/6} \frac{1+\cos(2\theta)}{2} \, d\theta \\
            &= 2 \int_{-\pi/6}^{\pi/6} (1+\cos(2\theta)) \, d\theta \\
            &= 2 \left[ \theta + \frac{\sin(2\theta)}{2} \right]_{-\pi/6}^{\pi/6} \\
            &= 2 \left[ \frac{\pi}{6} + \frac{\sin(\pi/3)}{2} - \left(-\frac{\pi}{6} + \frac{\sin(-\pi/3)}{2}\right) \right] \\
            &= 2 \left[ \frac{\pi}{6} + \frac{\sqrt{3}}{4} + \frac{\pi}{6} + \frac{\sqrt{3}}{4} \right] \\
            &= 2 \left[ \frac{\pi}{3} + \frac{\sqrt{3}}{2} \right] \\
            &= \frac{2\pi}{3} + \sqrt{3}
        \end{align*}
        \item The volume of the solid obtained by revolving the area under the curve about the $x$-axis can be calculated using the formula
        \begin{align*}
            V &= \pi \int_{-1}^{1} (\sqrt{4-x^2})^2 \, dx \\
            &= \pi \int_{-1}^{1} (4-x^2) \, dx \\
            &= \pi \left[ 4x - \frac{x^3}{3} \right]_{-1}^{1} \\
            &= \pi \left[ (4 - \frac{1}{3}) - (-4 + \frac{1}{3}) \right] \\
            &= \pi \left[ \frac{11}{3} + \frac{11}{3} \right] \\
            &= \frac{22\pi}{3}
        \end{align*}
        \item The centroid of the solid can be calculated using the formula
        \begin{align*}
            \bar{y} &= \frac{\frac{1}{2}\int_a^b [f(x)]^2 \, dx}{\int_a^b f(x) \, dx} \\
            &= \frac{\frac{1}{2}\int_{-1}^{1} (\sqrt{4-x^2})^2 \, dx}{\int_{-1}^{1} \sqrt{4-x^2} \, dx} \\
            &= \frac{\frac{1}{2}\int_{-1}^{1} (4-x^2) \, dx}{\frac{2\pi}{3} + \sqrt{3}} \\
            &= \frac{\frac{1}{2} \cdot \frac{22\pi}{3}}{\frac{2\pi}{3} + \sqrt{3}} \\
            &= \frac{11\pi}{3} \cdot \frac{1}{\frac{2\pi}{3} + \sqrt{3}} \\
            &= \frac{11\pi}{3} \cdot \frac{3}{2\pi + 3\sqrt{3}} \\
            &= \frac{11\pi}{2\pi + 3\sqrt{3}}
        \end{align*}
        Since the solid is symmetric about the $y$-axis, then $\bar{x}=0$. Thus, the centroid of the solid is at $\left(0, \frac{11\pi}{2\pi + 3\sqrt{3}}\right)$.
        \item The arc length of the curve in the first quadrant is one quarter of the circumference of a circle with radius 2, which can be calculated using the formula
        \begin{align*}
            K &= \frac{1}{4} \cdot 2\pi r = \frac{1}{4} \cdot 2\pi \cdot 2 = \pi
        \end{align*}
        \item The centroid of the surface area obtained by revolving the arc about the $y$-axes can be calculated using Pappus' theorem, which states that the centroid of a surface of revolution is located at a distance from the axis of revolution equal to the distance traveled by the centroid of the generating curve. Since the centroid of the arc is located at a distance of $\frac{4}{\pi}$ from the $y$-axis, then the centroid of the surface area is located at a distance of $\pi \cdot \frac{4}{\pi} = 4$ from the $y$-axis. Thus, the centroid of the surface area is at $\left(4, 0\right)$.
    \end{enumerate}
    \item The arc length of the curve can be calculated using the formula
    \begin{align*}
        K &= \int_{0}^{1} \sqrt{1 + \left(\frac{dy}{dx}\right)^2} \, dx \\
        &= \int_{0}^{1} \sqrt{1 + 25} \, dx \\
        &= \int_{0}^{1} \sqrt{26} \, dx \\
        &= \sqrt{26}
    \end{align*}
    \end{enumerate}
\end{document}


